% Reunión expositiva de precedencia causal entre resultados ya establecidos.
% Los propietarios íntegros continúan gobernando cada flecha.
% Propietarios reunidos por el grafo activo:
% - sucesor_102/deltas_ley9/c27_teorema_generacion_coinductiva_profundidad_arbitraria.tex
% - sucesor_102/deltas_ley9/c34_diagrama_bifurcado_y_decada_accion.tex
% - sucesor_102/deltas_ley9/c36_recta_accion_planck.tex
% - colaboracion/parte_iii/tex/capitulos_finales/32_raiz_menos_uno.tex
% - incorporaciones/parte_iv_ampliaciones_20260824/sources/editor_ready/03_pareja_planck_respuesta.tex
% - integracion_83/parte_iv/chapters/48_electron_primera_realizacion.tex

\section{Genealogía operatoria de la acción, la geometría angular y la masa}
\label{sec:l9s102:transicion-constantes-accion-geometria-masas}

La dinámica coinductiva APP--TRIT--TPK se prolonga a profundidad arbitraria
sobre un mismo estado enriquecido y sobre la estructura discreta conjunta del
continuo. En cada prolongación conserva valor, genealogía, fase, acarreo,
orientación, vacancia, ruta, frontera, memoria y supervivencia. Sus
publicaciones distinguidas son primero las tres coordenadas internas y,
mediante el sello dodecafásico, la coordenada \(\alpha_{\mathrm{HMT}}\):
\[
 \mathrm{APP}\longrightarrow\mathrm{TRIT}\longrightarrow\mathrm{TPK}
 \longrightarrow\mathcal S_{\infty}^{\mathrm{enr}}
 \longrightarrow\mathcal C_{\mathrm{cont}}^{\mathrm{disc}}
 \longrightarrow(\pi_{\mathrm{HMT}},\varphi_{\mathrm{HMT}},e_{\mathrm{HMT}})
 \longrightarrow\alpha_{\mathrm{HMT}}.
\]
La última flecha representa la publicación firmada de
\(\alpha_{\mathrm{HMT}}\) dentro de esa misma dinámica; no recibe valores
convencionales ni convierte a \(\pi,\varphi,e\) convencionales en generadores.

Sólo después se aplica el lector determinantal de acción. Éste publica la
valoración \(10^{-34}\mathcal U_S\), las dos secciones
\(\hbar_{\mathrm{pre}},\hbar_{\mathrm{ret}}\) y sus acciones de ciclo
completo
\[
 h_{\mathrm{pre}}=2\pi_{\mathrm{HMT}}\hbar_{\mathrm{pre}},
 \qquad
 h_{\mathrm{ret}}=2\pi_{\mathrm{HMT}}\hbar_{\mathrm{ret}}.
\]
A partir de este punto causal permanecen separadas dos ramas que no se
seleccionan entre sí. La rama dimensional--gravitatoria prolonga la acción y
la sección causal hasta las formas ya demostradas de la pareja de Planck,
\[
 \ell^{2}=G\hbar c^{-3},
 \qquad
 \frac{\ell}{Q}=Gc^{-4};
\]
la rama compleja publica, desde el operador interno de rotación,
\[
 u^2+1=0,
 \qquad u\longmapsto J_{+1},
 \qquad J_{+1}^{2}=-I.
\]
Ésta es la realización operatoria interna de \(i\), no la introducción de un
símbolo complejo sin antecedente. Las ramas dimensional--gravitatoria y
compleja son paralelas sobre el estado enriquecido común: la primera no
convierte a \(i\) en entrada de la escala de acción y, recíprocamente,
\(h\), \(\hbar\), \(c^{-3}\) y \(c^{-4}\) no generan a \(i\). Las identidades
de la pareja de Planck permanecen así separadas de la realización compleja.

Escribiendo \(\prec_{\mathrm{HMT}}\) para precedencia causal con conservación
de tipo ---y no para afirmar que cada objeto de la derecha sea función de
todos los de la izquierda---, se obtiene
\[
\begin{aligned}
 (\pi,\varphi,e,\alpha)_{\mathrm{HMT}}
 &\prec_{\mathrm{HMT}}(10^{-34}\mathcal U_S,\hbar,h)\\
 &\prec_{\mathrm{HMT}}
 \bigl\{G\hbar c^{-3},Gc^{-4};\,J_{+1}^{2}=-I, i\bigr\}\\
 &\prec_{\mathrm{HMT}}(A_{\mathrm{deg}},C^{*}_{\mathrm{deg}})\\
 &\prec_{\mathrm{HMT}}(\text{electrón},\text{ley de masas}).
\end{aligned}
\]
La construcción angular conserva sus entradas efectivas
\(\alpha_{\mathrm{HMT}}\) y \(\eta_{\mathrm{ret}}\): las dos ramas paralelas
precedentes fijan la disponibilidad conjunta y el tipado, pero no se promueven
retrospectivamente a selectores del ángulo o del contraángulo.

En la ruta electrónica central, los lectores bidireccionales conservan dos
etapas distintas de una sola genealogía:
\[
 \begin{aligned}
 E_e^{(0)}
   &=0.5109989224880660725\ldots\ \mathrm{MeV},
   &&\text{etapa basal},\\
 E_e^{\beta}
   &=E_e^{(0)}R_{\mathrm{act}}^{\,80-54\alpha_{\mathrm{HMT}}+6\Delta_4}
     =0.510998950690000808608\ldots\ \mathrm{MeV},
   &&\text{salida beta rectora},
 \end{aligned}
\]
con
\[
 K_D(\Gamma_e)=K_{\Omega}(\Gamma_e)=(80,54,6).
\]
La coincidencia del triple en la fibra central no identifica los dominios ni
las operaciones de \(K_D\) y \(K_{\Omega}\). En el registro canónico
histórico, \texttt{MD-ELECTRON-001} designa el candidato calibrado de la etapa
basal. Esa identificación conserva su función de procedencia, pero no gobierna
la salida beta ni selecciona la ruta; el basal y la clausura beta permanecen
como etapas tipadas de los lectores bidireccionales.

\paragraph{Criterio de refutación.}
La genealogía anterior queda refutada si un valor convencional o la salida
beta se usa para seleccionar una ruta, si la etapa basal se presenta como
cierre rector, si se identifica cualquiera de las dos ramas paralelas sin un
mapa tipado, o si el ángulo, el contraángulo, el electrón o una masa se colocan
como entradas del generador de constantes o de acción.
